\scanpage{228}
\mspage{174}
% archive sheet number: 227
enfin

\textcircled{c} \deleted{dans} \uline{$\mathrm{Fib}$}$_W\,S$, on a
\diagram{\begin{tikzcd}[ampersand replacement=\&,row sep=small,column sep=small]
W_S^f \arrow[r,Leftrightarrow] \arrow[d,Rightarrow] \& W_S^d \arrow[r,Leftrightarrow] \& W_S^g \arrow[r,Leftrightarrow] \& W_S^{\mathrm{univ}}\\
W_S \& {.} \& \&
\end{tikzcd}}

Cette implication $W_S^f\Rightarrow W_S$ est stricte si on ne fait pas
sur $W$ l'axiome $W(4\mathrm{b})$, \unsure{dont} l'utilité ici est
exactement d'\unsure{assurer}
\[
\left|\begin{aligned}
&W_S^f\Longleftrightarrow W_S\quad\text{dans }\underline{\mathrm{Fib}}_W\,S
  \quad\underline{\text{\uuline{si}}}\quad W(4\mathrm{b})\ \text{satisfaite}\\
&\text{et alors les cinq localiseurs envisagés ci-dessus}\\
&\qquad\text{dans }\underline{\mathrm{Fib}}_W\,S\text{ sont tous égaux.}
\end{aligned}\right.
\]

\deleted{\unsure{Mais}} Les relations \deleted{\unsure{entre}} localiseurs \deleted{\unsure{ci-dessus}} dans
$\underline{\mathrm{Liss}}_W\,S$, $\underline{\mathrm{Prop}}_W\,S$, $\underline{\mathrm{Fib}}_W\,S$ se transportent automatiquement aux sous-catégories de celles-ci, plus gén. à \unsure{toutes|toute} \unsure{celles|celle} qui s'envoient par un foncteur
d'\unsure{inclusion|induction} dans une de celles-ci. Dans le diagramme
p. \unsure{75|76}, ttes les catégories, à la seule exception
de $\mathrm{Cat}_{/S}$ tout entier, sont \deleted{dans} en ce cas.

\deleted{\unsure{\textcircled{d}}\illegible{1}}\inserted{Pour} les résultats principaux développés jusqu'à
présent sur $\mathrm{HOT}_W$, quand l'axiome $W(4\mathrm{b})$ est
utilisé, la forme faible \deleted{\unsure{(}}$W(4\mathrm{b}')$ suffit. \deleted{Celle-ci \unsure{assure}} On a
\[
\left|\begin{aligned}
&W_S^f\Longleftrightarrow W_S\quad\text{dans }\underline{\mathrm{Parf}}_W\,S
 \quad\underline{\text{\uuline{si}}}\quad W(4\mathrm{b}')\ \text{satisfaite}\\
&\text{et \inserted{alors} les cinq localiseurs }W_S^f,W_S^d,W_S^g,W_S^{\mathrm{univ}},W_S\\
&\text{sont égaux dans }\mathrm{Parf}_W\,S .
\end{aligned}\right.
\]

\textcircled{2} \uline{Principales catégories de catégories sur} $S$.

Ce sont, à des modifications d'éclairage près, celles
du diagramme p. \unsure{76}. Les catégories $\underline{\mathrm{Fibsc}}\,S$, $\underline{\mathrm{Fib}}\,S$,
$\underline{\mathrm{Cofibsc}}\,S$, $\underline{\mathrm{Cofib}}\,S$ et leurs variantes d'indice $0$,
sont inchangées -- leur rôle est d'être aussi petites
